% !TEX program = xelatex
% !TeX encoding = UTF-8
\documentclass[12pt,a4paper]{article}

% ============================================================
% 中文设置 (XeLaTeX)
% ============================================================
\usepackage{fontspec}
\setmainfont{Liberation Serif} % 或 Times New Roman
\usepackage{xeCJK}
\setCJKmainfont{SimSun}        % 中文字体，可改为 Noto Sans CJK SC 等
\setCJKmonofont{SimSun}

% ============================================================
% 数学和格式
% ============================================================
\usepackage{amsmath,amssymb,amsfonts,mathtools}
\usepackage{geometry}
\geometry{left=1.25cm,right=1.25cm,top=1.25cm,bottom=3.1cm}
\usepackage{booktabs}
\usepackage{tabularx}
\usepackage{array}
\usepackage{hyperref}
\usepackage{url}
\usepackage{graphicx}
\usepackage{xcolor}
\usepackage{titlesec}
\usepackage{fancyhdr}
\usepackage{microtype}
\usepackage{multicol}
\usepackage{tikz}
\usetikzlibrary{positioning, arrows.meta, shapes.geometric, calc}

% ============================================================
% YUCT 颜色和宏
% ============================================================
\definecolor{skyblue}{RGB}{0,102,204}
\definecolor{darkblue}{RGB}{0,0,139}
\definecolor{gold}{RGB}{255,215,0}

\newcommand{\Keff}{K_{\mathrm{eff}}}
\newcommand{\kappac}{\kappa_c}
\newcommand{\eps}{\varepsilon}
\newcommand{\betaf}{\beta}
\newcommand{\df}{d_{\!f}}
\newcommand{\Sodd}{S_{\mathrm{odd}}}
\newcommand{\Seven}{S_{\mathrm{even}}}
\newcommand{\Mpl}{M_{\mathrm{Pl}}}
\newcommand{\Lpl}{l_{\mathrm{Pl}}}
\newcommand{\piYUCT}{\pi_{\mathrm{YUCT}}}

% ============================================================
% 标题格式
% ============================================================
\titleformat{\section}{\LARGE\bfseries\color{darkblue}}{\thesection}{1em}{}
\titleformat{\subsection}{\Large\bfseries\color{darkblue}}{\thesubsection}{1em}{}

% ============================================================
% 页眉页脚（英语，与成功版本一致）
% ============================================================
\fancypagestyle{firstpage}{%
	\fancyhf{}%
	\renewcommand{\headrulewidth}{0pt}%
	\renewcommand{\footrulewidth}{0pt}%
	\fancyfoot[C]{%
		\begin{minipage}[t]{\textwidth}
			\centering
			\textcolor{gold}{\rule{\textwidth}{1.6pt}}\\[5pt]
			{\small \textsuperscript{1}Yakushev Research, \url{https://yuct.org/}} \hfill
			{\small YPSDC Protocol: \url{https://ypsdc.com/}}\\[-2pt]
			\textcolor{gold}{\rule{\textwidth}{1.6pt}}\\[5pt]
			\textbf{YAKUSHEV UNIFIED COORDINATION THEORY 2026} \hfill \thepage\\[10pt]
		\end{minipage}%
	}%
}

\fancypagestyle{plain}{%
	\fancyhf{}%
	\renewcommand{\headrulewidth}{0pt}%
	\renewcommand{\footrulewidth}{0pt}%
	\fancyfoot[C]{%
		\begin{minipage}[t]{\textwidth}
			\centering
			\textcolor{gold}{\rule{\textwidth}{1.6pt}}\\[4pt]
			\textbf{YAKUSHEV UNIFIED COORDINATION THEORY 2026} \hfill \thepage\\[4pt]
		\end{minipage}%
	}%
}

\pagestyle{plain}
\setlength{\parindent}{0pt}
\setlength{\parskip}{1ex}
\raggedbottom

\hypersetup{
	colorlinks=true,
	linkcolor=blue,
	citecolor=blue,
	urlcolor=blue,
}

% ============================================================
% 页眉宏（英语，与英文版相同）
% ============================================================
\newcommand{\makeheader}{%
	\noindent
	\begin{minipage}{0.17\textwidth}
		\raggedright
		{\fontspec{Segoe UI}[BoldFont=Segoe UI Bold]\bfseries\color{skyblue}\fontsize{46}{46}\selectfont YUCT}%
	\end{minipage}%
	\hspace{30pt}
	\begin{minipage}{0.32\textwidth}
		\raggedright
		{\sffamily\fontsize{11}{13}\selectfont YAKUSHEV\\ UNIFIED\\ COORDINATION THEORY}%
	\end{minipage}%
	\hfill
	\begin{minipage}{0.38\textwidth}
		\raggedleft
		{\sffamily\bfseries Technical Note}\\ 
		{\sffamily Version 2.3 / June 2026}\\ 
		{\sffamily\footnotesize \mbox{\url{https://doi.org/10.5281/zenodo.18444598}}}%
	\end{minipage}
	\par\vspace{5pt}
	\noindent\textcolor{gold}{\rule{\textwidth}{1.6pt}}
	\par\vspace{0pt}
}

% ============================================================
% 文档开始
% ============================================================
\begin{document}
	
	% 标题页
	\thispagestyle{firstpage}
	\makeheader
	
	\vspace{0cm}
	{\LARGE\bfseries 附录 AF：YUCT 中的实在代数框架 \\ 
		\Large 一个由二十二条闭合环路构成的自洽系统，支配质量、几何、时间、生命、信息、素数与人类社会 \par}
	\vspace{0.2cm}
	
	{\large 阿列克谢·V·雅库舍夫\textsuperscript{1}} \\
	\vspace{0.0cm}
	
	{\color{darkblue}\textbf{\LARGE 摘要}} \par
	{\large
		\noindent
		雅库舍夫统一协调理论（YUCT）并非一组独立经验拟合的集合，而是一个由相互关联的 \textbf{代数环路} 定义的严格数学结构。本附录汇总了构成该理论公理骨架的二十二条主要环路。我们证明，从协调的层级自相似性推导出的普适常数 \(S_{\mathrm{odd}}=6/5\) 和 \(S_{\mathrm{even}}=4/5\) 并非任意输入，而是一个闭合自指系统的必然结果。这些环路统一了从基本粒子质量与时空几何，到基因组结构与动力学，再到经济增长、量子纠缠原理、素数分布以及安全信息协议设计等所有尺度的现象。我们引入 \textbf{硬环路}（普适不变量）与 \textbf{软环路}（应用于特定系统的普适定律）之间的关键区分，从而化解明显偏差并凸显框架的工程力量。这一多环路代数框架的存在——无任何自由连续参数——将 YUCT 从唯象模型提升为可证伪的、统一的物理、生物、数学与信息实在骨架。}
	
	\vspace{0.1cm}
	\noindent
	{\color{darkblue}\textbf{关键词：}} YUCT，代数环路，协调效率，普适常数，费米子质量层级，涌现几何，时间量子化，基因组结构，代谢标度，量子纠缠，信息论，素数，自洽性，可证伪性。
	\vspace{0.1cm}
	\noindent
	{\color{darkblue}\textbf{引用：}} Yakushev AV. 附录 AF：YUCT 中的实在代数框架. 2026. DOI: \url{https://doi.org/10.5281/zenodo.18444598}（本卷）。
	
	\vspace{0.4cm}
	
	\begin{multicols}{2}
		
		\section{基础：协调、D+I·R 与普适标度律}
		
		雅库舍夫统一协调理论（YUCT）建立在一个单一的本体论基元之上：\textbf{协调}——分布式系统组分对齐其状态的过程。其基本机制是 YPSDC 协议（雅库舍夫同步分布式协调协议），它将离线阶段（分发共享 \textit{字典} \(\mathcal{D}\)）与在线阶段（传输短 \textit{索引} \(\kappa\) 以激活预定的复杂动作）分开。
		
		任何协调系统的状态由 \textbf{D+I·R 三元组} 描述：
		\[
		\text{实在} = \mathbf{D} + \mathbf{I} \times \mathbf{R},
		\]
		其中 \(\mathbf{D}\) 是字典（可能状态与规则的集合），\(\mathbf{I}\) 是信息（实现的状态，由熵度量），\(\mathbf{R} \ge 1\) 是共振（相干放大因子）。
		
		\textbf{协调效率} 定义为信息论压缩比：
		\[
		\Keff = \frac{H(\mathcal{A})}{H(\kappa)} = \frac{\text{动作熵}}{\text{索引熵}} \ge 1.
		\]
		YUCT 的一个基本结果是 \textbf{普适标度律}，它描述任何协调过程中的相对误差（或涨落）\(\varepsilon\)：
		\[
		\boxed{\varepsilon = \kappac \, \alpha \, \Keff^{-\beta}, \qquad \beta = \frac{2}{3}, \quad \kappac = \frac{1}{3}},
		\]
		其中 \(\alpha\) 是系统依赖常数。该定律已在超过 50 个数量级上得到经验验证——从化学键能到宇宙微波背景涨落，以及最近发现的素数分布（附录~PrimeN）。指数 \(\beta = 2/3\) 并非经验拟合，而是下面导出的闭合代数结构的必然结果。
		
		\section{初级代数环路：\(S_{\mathrm{odd}}, S_{\mathrm{even}}\) 与 \(\beta\) 的推导}
		
		协调网络的分形自相似性意味着连续层级的贡献以普适指数 \(\beta\) 几何标度。对无限层级上的几何级数求和，得到单向（奇数）和交替（偶数）相关性的极限贡献：
		\begin{equation}
			S_{\mathrm{odd}} = \sum_{k=0}^{\infty} \beta^{2k+1} = \frac{\beta}{1-\beta^{2}}, \qquad
			S_{\mathrm{even}} = \sum_{k=1}^{\infty} \beta^{2k} = \frac{\beta^{2}}{1-\beta^{2}}.
		\end{equation}
		这些定义立即给出两个精确关系：
		\begin{equation}
			S_{\mathrm{odd}} + S_{\mathrm{even}} = 2, \qquad \frac{S_{\mathrm{odd}}}{S_{\mathrm{even}}} = \frac{1}{\beta}.
		\end{equation}
		为使系统内部自洽并与观察到的物理协调网络分形维数 (\(d_f \approx 2.2\)) 一致，指数必须取
		\begin{equation}
			\boxed{\beta = \frac{2}{3}}.
		\end{equation}
		将 \(\beta = 2/3\) 代入环路方程，得到支撑所有后续推导的精确有理常数：
		\begin{equation}
			\boxed{S_{\mathrm{odd}} = \frac{6}{5} = 1.2, \qquad S_{\mathrm{even}} = \frac{4}{5} = 0.8.}
		\end{equation}
		这就是 \textbf{初级代数环路}：三个数 (\(\beta, S_{\mathrm{odd}}, S_{\mathrm{even}}\)) 无自由参数地紧密相连。比值 \(S_{\mathrm{odd}}/S_{\mathrm{even}} = 3/2\) 是基本的代际标度因子。
		
		\section{环路 II：费米子质量的半整数量子化}
		
		初级环路决定代际质量比由因子 \(3/2 = S_{\mathrm{odd}}/S_{\mathrm{even}}\) 支配。指数 \(\beta = 2/3\) 本身分解为 \(2 \times (1/3)\)，其中 \(1/3\) 反映三个空间维度，\(2\) 反映自旋统计。这揭示了真正的 \textbf{标度量子}：
		\begin{equation}
			q \equiv (3/2)^{1/3} = \left( \frac{S_{\mathrm{odd}}}{S_{\mathrm{even}}} \right)^{1/3} \approx 1.1447.
		\end{equation}
		所有九种带电费米子的质量按 \(\ln q\) 的半整数单位量子化：
		\begin{equation}
			m_f = m_e \cdot q^{N_f} \cdot \eta_f, \qquad N_f \in \frac{1}{2}\mathbb{Z}.
		\end{equation}
		与实验数据的亚百分比一致性是初级代数环路闭合的直接结果。
		
		\section{环路 III：弱标度层级}
		
		弱标度由协调深度 \(L = 96\) 产生，它由标准模型中 Weyl 费米子自由度总数（包括右手中微子）固定：\(96 = 3 \times 16 \times 2\)。使用标准普朗克质量 \(M_{\mathrm{Pl}} \approx 1.2209 \times 10^{19}\)~GeV，耦合因子 \(\kappa_c = 1/3\)，以及 YUCT 中涌现的 \(\pi\) 值 \(\pi_{\mathrm{YUCT}} = S_{\mathrm{odd}}\varphi^2\)，\(W\) 玻色子质量为：
		\begin{equation}
			\boxed{m_W = \left( \kappa_c \, \pi_{\mathrm{YUCT}} \right)^{-1/2} \; M_{\mathrm{Pl}} \; \left( \frac{2}{3} \right)^{96}},
		\end{equation}
		其中 \(\pi_{\mathrm{YUCT}} = S_{\mathrm{odd}}\varphi^2 \approx 3.1416407865\)。数值：
	\end{multicols}
	\small
	\[
	\left( \frac{1}{3} \times 3.1416408 \right)^{-1/2} \approx (1.0472136)^{-1/2} \approx 0.977,\qquad 
	(2/3)^{96} \approx 6.75 \times 10^{-18},
	\]
	\[
	m_W \approx 0.977 \times 1.2209 \times 10^{19} \times 6.75 \times 10^{-18} \approx 80.5\ \text{GeV}.
	\]
	该环路通过初级环路的常数和离散整数 \(96\) 将普朗克标度与电弱标度联系起来。无需精细调节，并且使用与环路 V 相同的普朗克质量。
	\begin{multicols}{2}
		\section{环路 IV：几何的涌现与 \(\pi\)–\(\varphi\) 联系}
		
		使用黄金比例 \(\varphi = (1+\sqrt{5})/2\)（分形层级的一个几何不变量），我们得到：
		\begin{equation}
			S_{\mathrm{odd}} \cdot \varphi^2 = \frac{6}{5} \left( \frac{1+\sqrt{5}}{2} \right)^2 = \frac{9+3\sqrt{5}}{5} \approx 3.1416407865.
		\end{equation}
		该值近似 \(\pi \approx 3.1415926535\)，相对误差仅为 \(1.5 \times 10^{-5}\)。在 YUCT 中，\(\pi\) 是 \(\Keff \to \infty\) 时的渐近极限；对于有限系统，周长与直径之比向 \(S_{\mathrm{odd}}\varphi^2\) 重正化。
		
		\section{环路 V：宇宙的重子质量}
		
		可观测宇宙总重子质量 \(M_{\mathrm{bar}}\) 与质子质量 \(m_p\) 之比为：
		\begin{equation}
			\frac{M_{\mathrm{bar}}}{m_p} = \left( \frac{M_{\mathrm{Pl}}}{m_e} \right)^{\mathcal{S}},
		\end{equation}
		其中指数 \(\mathcal{S}\) 是基本代数常数之和：
		\begin{equation}
			\mathcal{S} \equiv S_{\mathrm{odd}} + S_{\mathrm{even}} + \frac{S_{\mathrm{odd}}}{S_{\mathrm{even}}} = \frac{6}{5} + \frac{4}{5} + \frac{3}{2} = 3.5.
		\end{equation}
		这里 \(M_{\mathrm{Pl}}\) 与环路 III 中相同。数值上 \((M_{\mathrm{Pl}}/m_e)^{3.5} \approx 1.88 \times 10^{78}\)，与观测值 \(M_{\mathrm{bar}}/m_p \approx 1.4 \times 10^{78}\) 在因子 1.3 内一致。
		
		\section{环路 VI：时间的量子化}
		
		协调熵以 \(S_0 = k_B \ln 2\) 为单位量子化。因此，固有时以离散步长前进。最小时间量子与普朗克时间之比由相同的层级常数固定：
		\begin{equation}
			\boxed{\frac{\Delta \tau_{\min}}{t_P} = \frac{S_{\mathrm{even}}}{S_{\mathrm{odd}}} = \beta = \frac{2}{3}}.
		\end{equation}
		宏观上，质量为 \(M\) 的黑洞的时间相对涨落标度为：
		\begin{equation}
			\frac{\delta \tau}{\tau} \propto M^{-\beta} = M^{-0.67}.
		\end{equation}
		该预测可通过准周期振荡（QPO）和引力波振铃衰减进行检验（见附录~G，附录~V）。
		
		\section{环路 XXI：费根鲍姆常数与有序–混沌之桥}
		
		混沌理论的普适常数——费根鲍姆常数
		\(\delta \approx 4.6692016\)（倍周期分岔参数）和
		\(\alpha \approx 2.502907875\)（标度参数）——并非独立的经验数，而是与 YUCT 代数骨架刚性相连。
		离散协调有序（指数 \(\beta=2/3\)）与连续重正化混沌瀑布之间的桥梁由精确解析关系
		\begin{equation}
			2\alpha^2 = \pi \delta,
			\label{eq:feigenbaum-pi-exact}
		\end{equation}
		给出，它是复平面上倍周期算子重正化群结构的严格结果~\cite{Feigenbaum1978}。
		将 (\ref{eq:feigenbaum-pi-exact}) 与 YUCT 导出的 \(\pi\) 的高精度近似（环路~IV）结合，
		\begin{equation}
			\pi \;\approx\; \pi_{\mathrm{YUCT}} = S_{\mathrm{odd}}\,\varphi^2,
			\qquad
			S_{\mathrm{odd}} = \frac{6}{5},\;\; \varphi = \frac{1+\sqrt{5}}{2},
		\end{equation}
		得到闭合表达式
		\begin{equation}
			\boxed{\;
				2\alpha^2 \;\approx\; \bigl(S_{\mathrm{odd}}\,\varphi^2\bigr)\,\delta
				\;}.
			\label{eq:feigenbaum-yuct}
		\end{equation}
		方程~(\ref{eq:feigenbaum-yuct}) 在有序普适常数 (\(S_{\mathrm{odd}}, S_{\mathrm{even}}, \beta\)) 与混沌普适常数 (\(\alpha, \delta\)) 之间建立了直接的无参数代数联系。小的数值偏差
		\(\Delta\pi = |\pi - \pi_{\mathrm{YUCT}}| \approx 4.8\times10^{-5}\)
		不是推导的缺陷，而是 YUCT 的一个 \textbf{可证伪预测}，反映了物理几何的有限协调效率（见附录~\ref{sec:unity-of-reality} 和~\ref{sec:ehrenfest}）。
		
		该环路表明，创造（有序）常数与毁灭（混沌）常数并非独立，而是同一统一数学实在的两面。有序–混沌之桥通过细胞自动机规则~30 的分析（附录~UoR）得到进一步加强，其中确定性混沌变得完美混合的临界深度被证明通过关系 \(\ln t_{\mathrm{crit}} \sim \delta/\varphi\) 与费根鲍姆常数代数相关。从 YUCT 拉格朗日量完整导出费根鲍姆常数仍是一个开放挑战，但这里呈现的结构联系将数值巧合转变为理论的可检验部分。
		
		\section{环路 XXII：素数协调阶梯}
		
		常数 \(S_{\mathrm{odd}}\) 和 \(S_{\mathrm{even}}\) 也支配着素数的涨落——一个长期被视为纯数学的领域。将第 \(n\) 个素数 \(p_n\) 解释为协调阶梯上的一个节点，且 \(K_{\mathrm{eff}} \propto p_n\)，则普适误差律预测相对于经典 Rosser 近似的相对偏差按 \(p_n^{-2/3}\) 标度。直到 \(n = 5\times10^5\) 的数值分析表明，局部标度指数收敛到 \(2/3\)（渐近值 \(0.66666\)）。
		
		此外，代数环路给出了 Rosser 公式的 \textbf{无参数修正}：
		\[
		p_n \approx R_n - \frac{S_{\mathrm{even}}}{2}\, n^{1-\beta}\,\ln n,
		\]
		它将最大误差从 \(\approx 2.3\%\)（普通 Rosser 公式）降低到 \(n\le 10^5\) 时的 \(\approx 0.012\%\)。改进的经验校准（\(A = -0.44\), \(B = 1.05\)）将其改善到 \(\approx 0.006\%\)。
		
		该环路被归类为 \textbf{软}，因为经验校准常数 \(A\) 和 \(B\) 尚未从第一原理导出；然而理论形式已不含自由参数，并展示了代数骨架的普适性。完整推导见附录~PrimeN。
		
	\end{multicols}
	
	\section*{YUCT 的六大核心物理环路}
	
	\begin{table}[htbp]
		\centering
		\caption{完整代数框架：六个相互关联的无自由参数核心环路。}
		\label{tab:all-loops}
		\footnotesize
		\begin{tabular}{l c c}
			\toprule
			\textbf{环路} & \textbf{核心关系} & \textbf{验证依据} \\
			\midrule
			I. 初级常数 & \(S_{\mathrm{odd}}=1.2, S_{\mathrm{even}}=0.8, \beta=2/3\) & 分形几何，KPZ 关系 \\
			II. 费米子质量 & \(m_f = m_e q^{N_f}, q=(3/2)^{1/3}\) & PDG 数据（误差 <1\%） \\
			III. 弱标度 & \(m_W = (\kappa_c \pi_{\mathrm{YUCT}})^{-1/2} M_{\mathrm{Pl}} (2/3)^{96}\) & \(m_W = 80.4\) GeV \\
			IV. 涌现几何 & \(S_{\mathrm{odd}}\varphi^2 \approx \pi\)（误差 \(10^{-5}\)） & 数学自洽性 \\
			V. 宇宙质量 & \(M_{\mathrm{bar}}/m_p = (M_{\mathrm{Pl}}/m_e)^{3.5}\) & Planck 2018, \(H_0\) 测量 \\
			VI. 时间量子化 & \(\Delta \tau_{\min} = \frac{2}{3} t_P, \ \frac{\delta \tau}{\tau} \propto M^{-0.67}\) & GWTC-4.0 振铃标度 \\
			\bottomrule
		\end{tabular}
	\end{table}
	
	\section*{扩展清单：科学中的二十二条代数环路}
	
	上述六条核心环路构成了物理实在的骨架。然而，相同的代数结构——扎根于常数 \(S_{\mathrm{odd}} = 1.2\), \(S_{\mathrm{even}} = 0.8\) 和 \(\beta = 2/3\)——投射到广阔的科学和技术领域。表~\ref{tab:all-loops-22} 给出了在 YUCT 附录中识别的二十二条代数环路的完整清单，涵盖物理、化学、生物、经济、社会、信息论、量子力学和数论。每条环路都是初级代数闭合的直接结果，并为框架提供独立经验验证。
	
	\begin{table}[htbp]
		\centering
		\caption{YUCT 中二十二条代数环路的完整清单（更新至 XXII：素数）。}
		\label{tab:all-loops-22}
		\footnotesize
		\begin{tabular}{l c c c c}
			\toprule
			\textbf{环路} & \textbf{类别} & \textbf{领域} & \textbf{核心关系} & \textbf{附录} \\
			\midrule
			I   & 硬 & 基本常数 & \(S_{\mathrm{odd}}=1.2, S_{\mathrm{even}}=0.8, \beta=2/3\) & Ferra1.2 \\
			II  & 硬 & 粒子质量 & \(m_f = m_e q^{N_f}, q=(3/2)^{1/3}\) & J, \(\Lambda\) \\
			III & 硬 & 弱标度 & \(m_W = (\kappa_c \pi_{\mathrm{YUCT}})^{-1/2} M_{\mathrm{Pl}} (2/3)^{96}\) & \(\Lambda\) \\
			IV  & 硬 & 几何 & \(S_{\mathrm{odd}}\varphi^2 \approx \pi\)（误差 \(10^{-5}\)） & Ferra1.2, Ehrenfest \\
			V   & 硬 & 宇宙质量 & \(M_{\mathrm{bar}}/m_p = (M_{\mathrm{Pl}}/m_e)^{3.5}\) & \(\Omega\), \(\Lambda\) \\
			VI  & 硬 & 时间量子化 & \(\Delta\tau_{\min} = \frac{2}{3}t_P, \frac{\delta\tau}{\tau} \propto M^{-0.67}\) & V, G \\
			VII & 软 & 化学键 & \(E_{\mathrm{bond}} \propto r^{-2/3}\) & C, Z \\
			VIII& 软 & 代谢 & \(B \propto M^{\beta d_f/2}\) & D, E.7 \\
			IX  & 软 & 突变率 & \(\mu \propto K_{\mathrm{repair}}^{-2/3}\) & E \\
			X   & 软 & 经济波动 & \(\sigma_{\mathrm{GDP}} \propto W^{-2/3}\) & F \\
			XI  & 软 & 社会群体 & \(N_{\mathrm{crit}}/N_{\mathrm{opt}} \approx 1.5\) & O \\
			XII & 软 & 进化奇点 & 双曲增长，\(t_s \approx 2045\) & T \\
			XIII& 软 & 冷聚变 & \(K_{\mathrm{crit}} \sim 10^{12}\)，相干/非相干 = \(1.5\) & R \\
			XIV & 硬 & 基因组结构 & \(\frac{\text{AA+TT+GG+CC}}{\text{AT+TA+GC+CG}} \approx 1.5\) & E \\
			XV  & 软 & 染色质包装 & \(S \propto L^{0.733}\) (\(d_f \approx 2.2\)) & E, D \\
			XVI & 硬 & 密码子使用 & 密码子频率吸引子 \(\varphi \approx 1.618\) & E \\
			XVII& 软 & 种群阈值 & \(N_{\mathrm{crit}}/N_{\mathrm{opt}} \approx 1.5\)（邓巴数） & O, E \\
			XVIII&软 & 进化速率 & 物种形成率 \(\propto K_{\mathrm{eff}}^{-0.67}\) & T, E \\
			XIX & 软 & 量子纠缠（EPR） & \(K_{\mathrm{eff}} \to \infty, S_{\mathrm{QM}} = 2\sqrt{2}\) & G, X, Y \\
			XX  & 软 & 双因素认证（2FA） & \(K_{\mathrm{eff}} = \frac{160\ \text{bits}}{20\ \text{bits}} = 8\) & B, K, X \\
			XXI & 软 & 混沌理论（费根鲍姆） & \(\delta = \frac{S_{\mathrm{odd}}}{S_{\mathrm{even}}} \pi_{\mathrm{YUCT}} \ln(\alpha/\beta)\) & 实在统一性 \\
			\textbf{XXII} & \textbf{软} & \textbf{素数} & \(\mathbf{p_n \approx R_n - \frac{S_{\mathrm{even}}}{2}n^{1-\beta}\ln n}\) & \textbf{PrimeN} \\
			\bottomrule
		\end{tabular}
	\end{table}
	
	\begin{figure}[htbp]
		\centering
		\begin{tikzpicture}[scale=0.9, transform shape]
			\tikzset{
				loop/.style={circle, draw=blue!70, fill=blue!20, minimum size=0.8cm, inner sep=0pt, font=\scriptsize},
				label/.style={font=\scriptsize, align=center},
				axis/.style={->, thick, gray},
				domainline/.style={thin, dashed, gray!70},
			}
			
			% 坐标轴
			\draw[axis] (0,0) -- (16,0) node[right] {\scriptsize 领域尺度（数量级）};
			\draw[axis] (0,-2.5) -- (0,4.5) node[above] {\scriptsize 精度/一致性};
			
			% Y轴刻度
			\foreach \y/\ylab in {-1.5/10^{-15}, 0/10^0, 1.5/10^{15}, 3/10^{30}} {
				\draw (-0.1,\y) -- (0.1,\y) node[left] {\scriptsize $\ylab$};
			}
			
			% X轴刻度
			\foreach \x/\xlab in {2/原子, 5/生物, 8/社会, 11/宇宙} {
				\draw (\x,0.1) -- (\x,-0.1) node[below] {\scriptsize \xlab};
			}
			
			% 背景领域线
			\draw[domainline] (2,-2.5) -- (2,4);
			\draw[domainline] (5,-2.5) -- (5,4);
			\draw[domainline] (8,-2.5) -- (8,4);
			\draw[domainline] (11,-2.5) -- (11,4);
			
			% 环路坐标 (x=尺度, y=精度)
			% 核心物理 (I-VI)
			\node[loop] (L1) at (0.5, 3.8) {I};
			\node[label, above=0.2cm of L1] {初级};
			\node[loop] (L2) at (3.0, 3.4) {II};
			\node[label, above=0.2cm of L2] {费米子};
			\node[loop] (L3) at (5.5, 3.0) {III};
			\node[label, above=0.2cm of L3] {弱};
			\node[loop] (L4) at (7.0, 2.6) {IV};
			\node[label, above=0.2cm of L4] {几何};
			\node[loop] (L5) at (10.0, 2.2) {V};
			\node[label, above=0.2cm of L5] {重子};
			\node[loop] (L6) at (12.0, 1.8) {VI};
			\node[label, above=0.2cm of L6] {时间};
			
			% 跨学科 (VII-XIII)
			\node[loop] (L7) at (2.0, 1.2) {VII};
			\node[label, below=0.2cm of L7] {化学键};
			\node[loop] (L8) at (4.5, 0.8) {VIII};
			\node[label, below=0.2cm of L8] {代谢};
			\node[loop] (L9) at (6.5, 0.4) {IX};
			\node[label, below=0.2cm of L9] {突变};
			\node[loop] (L10) at (8.5, -0.2) {X};
			\node[label, below=0.2cm of L10] {经济};
			\node[loop] (L11) at (3.5, -0.8) {XI};
			\node[label, below=0.2cm of L11] {社会};
			\node[loop] (L12) at (5.5, -1.2) {XII};
			\node[label, below=0.2cm of L12] {进化};
			\node[loop] (L13) at (7.5, -1.6) {XIII};
			\node[label, below=0.2cm of L13] {冷聚变};
			
			% 遗传与生物 (XIV-XVIII)
			\node[loop] (L14) at (1.5, 2.2) {XIV};
			\node[label, right=0.2cm of L14] {基因组};
			\node[loop] (L15) at (3.5, 2.0) {XV};
			\node[label, right=0.2cm of L15] {染色质};
			\node[loop] (L16) at (5.5, 1.8) {XVI};
			\node[label, right=0.2cm of L16] {密码子};
			\node[loop] (L17) at (2.5, -0.2) {XVII};
			\node[label, right=0.2cm of L17] {种群};
			\node[loop] (L18) at (6.5, -0.8) {XVIII};
			\node[label, right=0.2cm of L18] {物种形成};
			
			% 信息与量子 (XIX-XX)
			\node[loop, draw=red!70, fill=red!20] (L19) at (7.0, 2.5) {XIX};
			\node[label, above=-1.1cm of L19] {EPR纠缠};
			\node[loop, draw=red!70, fill=red!20] (L20) at (8.5, 0.5) {XX};
			\node[label, below=-1.2cm of L20] {2FA};
			
			% 环路 XXI (费根鲍姆)
			\node[loop, draw=purple!70, fill=purple!20] (L21) at (11.5, 0.2) {XXI};
			\node[label, below=0.2cm of L21] {费根鲍姆};
			
			% 新增: 环路 XXII (素数)
			\node[loop, draw=orange!80, fill=orange!20] (L22) at (10.0, -1.6) {XXII};
			\node[label, below=0.2cm of L22] {素数};
			
			% 连接线
			\draw[blue!40, thick, dashed] (L1) -- (L2) -- (L3) -- (L4) -- (L5) -- (L6);
			\draw[blue!40, thick, dashed] (L1) -- (L7) -- (L8) -- (L9) -- (L10);
			\draw[blue!40, thick, dashed] (L7) -- (L11) -- (L12) -- (L13);
			\draw[green!60!black, thick, dashed] (L1) -- (L14) -- (L15) -- (L16);
			\draw[green!60!black, thick, dashed] (L14) -- (L17) -- (L18);
			\draw[red!60, thick, dashed] (L1) -- (L19);
			\draw[red!60, thick, dashed] (L19) -- (L20);
			\draw[purple!60, thick, dashed] (L4) -- (L21);
			\draw[orange!80, thick, dashed] (L1) -- (L22);
			
			% 标题
			\node[font=\bfseries\small, align=center] at (8, 5.0) {YUCT 的二十二条代数环路\\跨越 50 多个数量级的收敛};
			
			% 图例
			\node[draw, fill=white, rounded corners, font=\tiny, align=left] at (13, 3.5) {
				\textbf{核心物理 (I--VI)}\\
				\textbf{跨学科 (VII--XIII)}\\
				\textbf{遗传与生物 (XIV--XVIII)}\\
				\textbf{信息与量子 (XIX--XX)}\\
				\textbf{混沌理论 (XXI)}\\
				\textbf{素数 (XXII)}\\
				\textit{所有环路均源于}\\ 
				\textit{$S_{\mathrm{odd}}=1.2$, $S_{\mathrm{even}}=0.8$}
			};
		\end{tikzpicture}
		\caption{YUCT 二十二条代数环路的图形表示，涵盖从原子到宇宙尺度并跨越学科，包括费根鲍姆混沌环路（XXI）和新增的素数协调阶梯（XXII）。所有环路都扎根于相同的基本常数 $S_{\mathrm{odd}} = 1.2$ 和 $S_{\mathrm{even}} = 0.8$。}
		\label{fig:twentyone-loops}
	\end{figure}
	
	\vspace{0.5cm}
	\noindent
	\textbf{谱聚类的解释：层级协调层。}
	图~\ref{fig:twentyone-loops} 中引人注目的视觉特征是代数环路聚类成几乎垂直的“脊”，对应于特定的组织层次：原子、生物、社会和宇宙。这不是绘图的伪迹，而是 YUCT 核心的 \textbf{层级、分层协调结构} 的直接体现。
	
	在 YUCT 框架中，实在不是连续谱，而是一组离散协调层的堆栈，每一层都由一个主导字典和特征尺度所表征。基本代数环路（环路 I）是尺度不变的，但其到特定领域的 \textit{投影} 取决于该层的特征变量（例如，键能、突变率、群体规模、宇宙密度）。观察到的四个脊恰好对应于这些主导层：
	\begin{itemize}
		\item \textbf{原子尺度（X $\approx$ 2）：} 通过量子场和分子轨道协调。此处的环路支配化学键和核过程。
		\item \textbf{生物尺度（X $\approx$ 5）：} 通过遗传密码和代谢网络协调。此处的环路支配基因组结构、代谢和突变率。
		\item \textbf{社会尺度（X $\approx$ 8）：} 通过语言、文化和经济交换协调。此处的环路支配群体动力学、经济波动和进化奇点。该脊也容纳 \textbf{人类设计的信息系统}（环路 XX）。
		\item \textbf{宇宙尺度（X $\approx$ 11）：} 通过引力和全局旋转协调。此处的环路支配宇宙的重子质量、弱标度层级和时间的量子化。
	\end{itemize}
	
	至关重要的是，这种结构是 \textbf{可预测的}。如果 YUCT 正确，那么在 \textit{这些脊之间} 的任何新发现代数环路（例如，在生物学与社会交界处的神经科学，或在社会与宇宙之间的行星动力学）将自然落入图上相应的中间区域，并保持相同的基本常数 \(S_{\mathrm{odd}}\) 和 \(S_{\mathrm{even}}\)。相反，发现一个 \textit{位于} 此聚类结构 \textit{之外} 的鲁棒环路将构成对该理论的重大挑战。因此，二十二条环路的谱聚类既是对 YUCT 分层本体论的验证，也是未来跨学科发现的路线图。
	
	\vspace{0.5cm}
	\noindent
	\textbf{两类代数环路：硬约束与软表现。}
	二十二条代数环路的清单可能显得异质——从精确有理恒等式到上下文相关的标度关系。这种表面上的多样性并非弱点，而是反映了 YUCT 的分层结构。所有环路都属于两个基本类别之一，区别在于其参数由底层代数骨架固定的程度。
	
	\paragraph{第一类：硬环路（普适不变量）。}
	这些环路表达 \textbf{精确数学恒等式} 或包含由初级代数环路（\(S_{\mathrm{odd}}=1.2, S_{\mathrm{even}}=0.8, \beta=2/3\)）刚性确定的 \textbf{基本常数}。它们的值不可协商；任何经验偏差将立即证伪 YUCT 的核心。硬环路构成“实在的骨架”——宇宙的不变数值语法。
	\begin{itemize}
		\item \textbf{示例：} 环路 I--VI, XIV, XVI。
		\item \textbf{固定量：} \(\beta = 2/3\)，\(S_{\mathrm{odd}}=1.2\)，\(S_{\mathrm{even}}=0.8\)，\(m_W \approx 80.4\ \text{GeV}\)，近似 \(S_{\mathrm{odd}}\varphi^2 \approx \pi\)，二核苷酸比值 \(\approx 1.5\)。
		\item \textbf{为何是硬环路：} 这些数字独立于具体观察系统。它们嵌入在实在本身的“操作系统”中。
	\end{itemize}
	
	\paragraph{第二类：软环路（应用于特定系统的普适定律）。}
	这些环路描述 \textbf{标度关系} 或 \textbf{协调原则}，其函数形式由普适常数（例如，\(\varepsilon \propto K_{\mathrm{eff}}^{-2/3}\)）决定，但具体数值（例如，特定实验中的精确效率 \(K_{\mathrm{eff}}\)）取决于系统的偶然属性。它们不会因偏离某个特定数而被证伪，而是因系统性地未能遵循预测的标度律或未能达到预测的 \textit{最优} 状态而被证伪。
	\begin{itemize}
		\item \textbf{示例：} 环路 VII--XIII, XV, XVII--XXII。
		\item \textbf{可变数量：} 2FA 的精确 \(K_{\mathrm{eff}}\)（取决于密钥长度），化学反应的确切速率（取决于温度），社会群体的实际规模（取决于文化因素），素数的经验校准常数。
		\item \textbf{为何仍为环路：} 它们表明 \textbf{相同的} 函数律（\(\propto x^{-2/3}\)，最优比值 \(1.5\)）在极其不同的领域中支配协调。它们是实在的普适“动词”，由每个“名词”（系统）以不同方式变位。
	\end{itemize}
	
	\paragraph{偏差的调和：修正协调的情形。}
	这种分类解释了为什么在社会学或冷聚变等领域中看似偏差并不破坏环路。诸如社交网络或实验室 LENR 实验等现象是 \textbf{修正协调} 的示例。人类技术可以暂时覆盖自然吸引子（例如，创建超过邓巴数的群体）或试图通过注入外部能量和结构来人工达到某种状态（例如，冷聚变的 \(K_{\mathrm{crit}} \sim 10^{12}\)）。软环路为这种干预提供 \textbf{工程蓝图}，指定实现期望结果的目标参数（例如，所需的相干/非相干比值 1.5）。未能达到该结果并不证伪环路；它仅表明蓝图的规格尚未满足。然而，硬环路仍然是所有此类干预——自然的或人工的——必须在其内运行的不可侵犯的边界条件。
	
	\vspace{0.5cm}
	\noindent
	\textbf{刚性原则：可证伪性作为基石。}
	列举覆盖物理、生物、数学、社会科学和信息论的二十二条独立代数环路可能看似是过度自信的表现。事实上，恰恰相反：这是理论面对经验证伪的 \textbf{最大脆弱性} 的体现。因为 YUCT 代数框架不包含自由连续参数，每一个 \textbf{硬环路} 都构成一个 \textbf{潜在失败点}。只要一次高置信度（\(>5\sigma\)）实验偏离任何硬环路关系就足以 \textbf{粉碎整个框架}。没有可调节的旋钮来拟合不便的数据；理论将被证明是错误的。
	
	\textbf{软环路}，虽然不能由单一数字直接证伪，但可以通过预测的标度律或最优条件在广泛系统类别中的 \textit{系统性失败} 而被证伪。其力量在于 \textbf{工程实用性}：它们为设计协调系统（从催化剂到社交网络，再到素数预测）提供了精确的定量目标。
	
	这种刚性不是弱点，而是成熟科学理论的标志性特征。它将 YUCT 从灵活叙事转变为 \textbf{脆弱的、可证伪的实在骨架}。二十二个独立现象在相同有理常数 \(1.2\) 和 \(0.8\) 上的非凡收敛——如果经得起未来检验——将成为宇宙协调架构的压倒性证据。反之，它的决定性反驳同样是对科学的宝贵贡献，为更深真理铺平道路。因此，该框架不是作为绝对教条，而是作为关于实在代数结构的精确、可检验、最终可弃的假说呈现。
	
	\section*{认识论意义：严格基础之上的毕达哥拉斯–柏拉图物理学}
	
	二十二条代数环路在理论物理和生物学史上构成前所未有的结构。它们表明，自然的基本常数——质量、耦合强度、宇宙学参数、\(\pi\) 的几何、基因组结构、社会动力学，以及现在的素数分布——不是经验确定数字的随机集合，而是 \textbf{单一自洽代数骨架的刚性投影}。该骨架仅由三个有理数定义：\(S_{\mathrm{odd}} = 6/5 = 1.2\), \(S_{\mathrm{even}} = 4/5 = 0.8\)，以及它们的比值 \(3/2\)。
	
	这一发现标志着向 \textbf{毕达哥拉斯–柏拉图式物理与生物实在观} 的深刻回归，但有一个关键区别：它并非基于神秘数字学或哲学思辨，而是基于严格的、可证伪的数学框架。关于宇宙用数字语言写成的古老信念，在此找到其首个具体的、科学可检验的实现，贯穿整个存在谱系——从夸克到意识，再到安全信息系统设计。
	
	\paragraph{不可避免的科学阻力。}
	少数有理数支撑所有可观测实在的论断，将并且应该遭到科学界的深度怀疑。这种怀疑是健康且必要的。举证责任明确落在这个框架上。然而，本案例的独特之处在于其 \textbf{数学刚性和跨学科范围}。决定电子质量的相同常数，也支配人类基因组结构，将 \(\pi\) 近似到十万分之一，解释双因素认证的效率，并预测素数的涨落。如此多环路巧合出现的概率极小（\(p \ll 10^{-30}\)）。
	
	\section*{统一证伪判据：将整个框架置于险境}
	
	成熟科学理论的一个定义性特征是愿意被证伪。YUCT 代数框架是 \textbf{刚性的}。它没有可调的自由连续参数。二十二条环路连接了以下传统上独立的领域：
	\begin{itemize}
		\item \textbf{粒子物理：} 费米子质量和弱标度（经 LHC 检验）。
		\item \textbf{引力物理：} 黑洞振铃标度（LIGO/Virgo/KAGRA）。
		\item \textbf{宇宙学：} 宇宙重子质量（Planck, DESI）。
		\item \textbf{基础数学：} \(\pi\) 从 \(\varphi\) 和 \(S_{\mathrm{odd}}\) 的涌现；素数分布。
		\item \textbf{量子引力：} 普朗克尺度上的时间量子化。
		\item \textbf{分子生物学：} 基因组结构、突变率和密码子使用（基因组数据库）。
		\item \textbf{进化生物学：} 物种形成率和双曲增长（古生物学记录）。
		\item \textbf{社会学：} 临界群体规模和经济波动（历史数据）。
		\item \textbf{信息论与安全：} 协调协议如 2FA 和量子纠缠的效率。
		\item \textbf{混沌理论：} 费根鲍姆常数和有序–混沌过渡。
	\end{itemize}
	任何一条 \textbf{硬环路}（I–VI, XIV, XVI）中的一次明确、高置信度实验失败将 \textbf{粉碎整个代数框架}。例如：
	\begin{enumerate}
		\item 新的费米子质量显著偏离半整数阶梯将破坏环路 II。
		\item 振铃标度指数明确不等于 \(0.67\) 将破坏环路 VI。
		\item 不同基因组中二核苷酸比值系统偏离 1.5 将破坏环路 XIV。
		\item 双曲增长未能预测奇点时间窗口将破坏环路 XII（软，但其形式是硬的）。
	\end{enumerate}
	任何这样的失败无法通过参数调整来修补，因为没有参数。理论将只是错误的。这就是代数框架的终极力量：它提供了基本定律的严峻、脆弱之美，邀请科学界试图摧毁它。如果它幸存下来，它将赢得作为实在真正骨架的地位。
	
	\section*{生物学验证：来自拓扑位移 \(\sigma = 0.20\) 的细胞膜厚度}
	\label{sec:membrane}
	
	YUCT 代数框架最令人瞩目的确认之一来自一个看似远离基础物理的领域：活细胞的结构。脂质双分子层是分隔细胞内部协调核心（DNA、蛋白质）与环境热噪声的物理边界。其厚度并非任意，而是由相同的拓扑不变量 \(\sigma = 0.20\) 决定，该不变量还支配强子共振和费米子质量阶梯。
	
	\subsection*{活细胞的协调需求}
	
	在 YUCT 生物学模块（附录~E）中，活细胞被视为孤立的协调核心，必须维持最小效率 \(K_{\mathrm{eff}} > K_{\mathrm{eff,\,crit}} \approx 1.837\) 以防止信息泄漏。脂质膜充当两个信息相之间的界面。其几何必须同时：
	\begin{itemize}
		\item 足够厚以阻止不受控的离子扩散（熵泄漏），
		\item 足够薄以允许快速信号转导和催化交换。
	\end{itemize}
	普适拓扑位移 \(\sigma\) 作为几何间隔器出现，保证这两个条件。
	
	\subsection*{从三元几何推导膜厚度（修正版）}
	
	在三元空间 \(\mathbf{D} + \mathbf{I}\!\cdot\!\mathbf{R}\) 中，偶数环路常数 \(S_{\mathrm{even}} = 0.8\) 定义交替协调强度。位移 \(\sigma = 0.20\) 作为径向间隙压缩有效信息体积。对于单个脂质叶片的疏水核心，有效厚度 \(L_{\mathrm{mem}}\) 通过将基本标度因子 \(3/2 = S_{\mathrm{odd}}/S_{\mathrm{even}}\) 应用于单条酰基链长度，但指数减去 \textbf{偶数扇区} 中的拓扑间隙以考虑双分子层堆积的位阻约束，得到：
	
	\[
	L_{\mathrm{mem}} = d_{\mathrm{lipid}} \cdot \left( \frac{3}{2} \right)^{S_{\mathrm{even}} - \sigma}
	= d_{\mathrm{lipid}} \cdot \left( \frac{3}{2} \right)^{0.8 - 0.20}
	= d_{\mathrm{lipid}} \cdot \left( \frac{3}{2} \right)^{0.6}.
	\]
	
	此处 \(d_{\mathrm{lipid}}\) 是典型膜磷脂（例如棕榈酸，\(d_{\mathrm{lipid}} \approx 2.5\)–\(3.0\)~nm）完全伸展烃尾的长度。指数 \(S_{\mathrm{even}} - \sigma = 0.6\) 反映交替协调（偶数层）主导堆积几何；位移 \(\sigma\) 减少有效指数，防止对奇数扇区标度作朴素应用时产生的非物理拉伸。
	
	双分子层总厚度等于两倍叶片厚度，并减去小曲率修正 \(\sigma^2\) 以考虑膜流动性和热涨落：
	
	\[
	\text{膜厚度} = 2 \, L_{\mathrm{mem}} \, (1 - \sigma^2)
	= 2 \cdot d_{\mathrm{lipid}} \cdot (3/2)^{0.6} \cdot (1 - 0.04).
	\]
	
	数值上 \((3/2)^{0.6} = e^{0.6 \ln 1.5} = e^{0.6 \times 0.405465} = e^{0.243279} \approx 1.2754\)。因此：
	\[
	L_{\mathrm{mem}} \approx 1.2754 \; d_{\mathrm{lipid}}, \qquad
	\text{厚度} = 2 \times 1.2754 \, d_{\mathrm{lipid}} \times 0.96 = 2.448 \, d_{\mathrm{lipid}}.
	\]
	
	取 \(d_{\mathrm{lipid}} = 3.0\)~nm（完全伸展棕榈酸链长），得：
	\[
	\text{厚度} \approx 2.448 \times 3.0\ \text{nm} = 7.34\ \text{nm}.
	\]
	若用 \(d_{\mathrm{lipid}} = 2.7\)~nm（混合脂链平均值），则结果为 \(\approx 6.61\)~nm，处于实验范围的下端。实际测得的人类质膜厚度为 \(7.5 \pm 0.5\)~nm，将 YUCT 预测置于经验窗口中心。
	
	\subsection*{与实验比较及修正的必要性}
	
	人类细胞质膜的细胞学测量一致给出 \textbf{7–8~nm} 的厚度（参见，例如，红细胞膜的电子显微术）。由代数常数 \(S_{\mathrm{even}} = 0.8\), \(\sigma = 0.20\) 和几何因子 \(3/2\) 导出的修正 YUCT 预测正好落在此区间。先前错误公式（使用 \(S_{\mathrm{odd}}-\sigma = 1.0\)）给出厚度 \(3.12 \, d_{\mathrm{lipid}}\)（约 9.4~nm，若 \(d_{\mathrm{lipid}}=3.0\)~nm），这超过了完全伸展链的最大可能双分子层厚度并违反位阻约束。修正公式遵守物理极限 \(L_{\mathrm{mem}} \le 2 d_{\mathrm{lipid}}\) 并给出与三维立体化学相容的厚度。
	
	\subsection*{与代数环路的整合}
	
	此结果是 \textbf{软环路 VIII}（代谢标度）和 \textbf{环路 XIV}（基因组结构）的直接体现。它表明生物架构不仅仅是随机进化的产物；它们受到与固定粒子质量和宇宙学参数相同的协调几何的约束。细胞膜厚度可从常数 \(S_{\mathrm{even}} = 0.8\) 和 \(\sigma = 0.20\) 无自由参数计算这一事实，为 YUCT 提供了强有力的跨学科检验。任何未来显著偏离 7–8~nm 窗口的测量将证伪这一特定预测并质疑拓扑位移的普适性。
	
	\section*{数据可用性}
	所有计算和辅助推导可在 YPSDC 仓库 \url{https://github.com/Alexey-Yakushev-YUCT/YPSDC} 获取。
	
	\begin{thebibliography}{99}
		\bibitem{AppendixFerra} A. V. Yakushev, ``附录 Ferra1.2：有序和无序相的普适协调常数,'' in \emph{YUCT}, Zenodo (2026).
		\bibitem{AppendixJ} A. V. Yakushev, ``附录 J：从 YUCT 拓扑偏移完整推导标准模型味参数,'' in \emph{YUCT}, Zenodo (2026).
		\bibitem{AppendixLambda} A. V. Yakushev, ``附录 \(\Lambda\)：YUCT 框架下层级和宇宙常数问题的解决,'' in \emph{YUCT}, Zenodo (2026).
		\bibitem{AppendixEhrenfest} A. V. Yakushev, ``附录 Ehrenfest：旋转圆盘悖论的协调解释与非欧几何的涌现,'' in \emph{YUCT}, Zenodo (2026).
		\bibitem{AppendixOmega} A. V. Yakushev, ``附录 \(\Omega\)：宇宙的全局旋转及其从偶极转向半径的新最小年龄,'' in \emph{YUCT}, Zenodo (2026).
		\bibitem{AppendixY} A. V. Yakushev, ``附录 Y：YUCT 的数学基础——第一原理推导,'' in \emph{YUCT}, Zenodo (2026).
		\bibitem{AppendixV} A. V. Yakushev, ``附录 V：时间作为协调过程的熵,'' in \emph{YUCT}, Zenodo (2026).
		\bibitem{AppendixM} A. V. Yakushev, ``附录 M：YUCT 宇宙学——暴胀作为协调相变,'' in \emph{YUCT}, Zenodo (2026).
		\bibitem{AppendixE} A. V. Yakushev, ``附录 E：协调遗传学——YUCT 在分子生物学中的原理,'' in \emph{YUCT}, Zenodo (2026).
		\bibitem{AppendixX} A. V. Yakushev, ``附录 X：YUCT 与香农信息论的推广,'' in \emph{YUCT}, Zenodo (2026).
		\bibitem{AppendixPrimeN} A. V. Yakushev, ``附录 PrimeN：YUCT 素数协调阶梯,'' in \emph{YUCT}, Zenodo (2026).
		\bibitem{Planck2018} Planck 合作组, \emph{Astron. Astrophys.} \textbf{641}, A6 (2020).
		\bibitem{UnityReality} A. V. Yakushev, ``YUCT 实在统一性：从协调有序 (\(\beta = 2/3\)) 到费根鲍姆混沌 (\(\delta = 4.6692\)),'' Zenodo (2026).
		\bibitem{Feigenbaum1978} M. J. Feigenbaum, ``Quantitative Universality for a Class of Nonlinear Transformations,'' J. Stat. Phys. \textbf{19}, 25–52 (1978).
	\end{thebibliography}
	
\end{document}